% source: manti.01.ma, 1ère SM, série « généralités sur les fonctions » (Drive 0B6pfSJcEVng0LWVnZDJaQUxRLU0).
% p.1 Ex.1-5 = série française fonctions_s4 ; p.2 = série française fonctions_s5 -> non repris
\ex{Fonctions}{LWVn}{Ex. 6}{11}
On considère la fonction $f$ définie par $f(x)=\dfrac{\sqrt{2x-1}}{x}$.
\begin{enumerate}
\item Déterminer $D_f$ et montrer que la fonction $f$ est minorée.
\item Montrer que $f$ est majorée par $1$.
\item On pose $g(x)=\dfrac{2x}{x^2+1}$ et $h(x)=\sqrt{2x-1}$.
\begin{enumerate}
\item Montrer que $T_g(x,y)=\dfrac{2(1-xy)}{\left(x^2+1\right)\left(y^2+1\right)}$ \ (l'original écrit $\dfrac{1-xy}{\dots}$ [?]).
\item Étudier la monotonie de $g$ sur $[1,+\infty[$ et sur $[0,1]$.
\item Vérifier que $f=g\circ h$, puis étudier la monotonie de $f$.
\end{enumerate}
\end{enumerate}
\fin
\ex{Fonctions}{LWVn}{Ex. 7}{7}
\begin{enumerate}
\item Soit $n$ un entier naturel non nul. Déterminer :
$E\!\left(\dfrac{n+1}{n}\right)$, \ $E\!\left(\dfrac{2n+3}{n+1}\right)$, \ $E\!\left(\sqrt{n^2+n}\right)$, \ $E\!\left(\sqrt{4n^2+4n}\right)$, \ $E\!\left(\sqrt{(n+1)(n+3)}\right)$.
\item Résoudre dans $\mathbb{R}$ les équations : $E(1-2x)=3$ ; \ $E\!\left(x^2-3x+3\right)=1$ ; \ $E\!\left(\sqrt{x-2}\right)-2=0$.
\item Résoudre dans $\mathbb{R}$ les inéquations : $E(3-x)<1$ ; \ $E(x)\geq2$ ; \ $2E\!\left(\dfrac{2}{x-1}\right)\leq3$.
\end{enumerate}
\fin
\ex{Fonctions}{LWVn}{Ex. 8}{7}
On considère la fonction $f$ définie par $f(x)=E(2x)-2x$.
\begin{enumerate}
\item Montrer que $f$ est périodique de période $T=\dfrac12$.
\item Écrire l'expression de $f(x)$ sur chacun des intervalles $[0,1[$ et $[1,2[$.
\item Tracer la courbe $(C_f)$ sur l'intervalle $[-3,5[$.
\end{enumerate}
\fin
